\begin{frame}{softmax: 원점수를 확률로 --- 숫자로}
  \begin{itemize}
    \item 모델이 산출하는 것은 확률이 아니라 \textbf{원점수}
      (logit) $z_j \in \mathbb{R}$. \kb{softmax}(Week 2.3
      시그모이드의 다클래스 일반화):
      \[P(x_t = j) = \frac{e^{z_j}}{\sum_{k} e^{z_k}}\]
    \item 왜 $e^{z}$로 정규화: (1) 모든 확률 0$\sim$1, 합 1,
      (2) \textbf{원점수 차이가 그대로 확률 비율로}:
      $z_a - z_b = 1$이면 $P(a)/P(b) = e \approx 2.7$배.
    \item 숫자로: $z = (4, 1, 0, -3)$ 이면
      $e^z \approx (54.6, 2.72, 1.0, 0.05)$, 합 $\approx 58.37$
      $\to$ 확률 $\approx$ \textbf{(0.935, 0.047, 0.017, 0.001)}
      — 첫 토큰에 93.5\%가 모인 ``확신 있는 예측''.
    \item 반면 $z = (1,1,1,1)$이면 (0.25, 0.25, 0.25, 0.25) —
      \textbf{완전한 무지(무차별 우연)}.
    \item \textbf{shift-invariance}: 모든 logit에 같은 상수를
      더해도 분포는 변하지 않음 — 확률은 \textbf{차이}로
      결정. 이 벡터를 실제 정답과 비교하는 것이 교차
      엔트로피 손실.
  \end{itemize}
\end{frame}
